% generated by figures/make_tables.py; do not edit by hand
\begingroup\scriptsize
\begin{longtable}{@{}p{0.20\textwidth} p{0.09\textwidth} p{0.16\textwidth} p{0.04\textwidth} p{0.07\textwidth} p{0.30\textwidth}@{}}
\caption{Fields of the \texttt{model-metrics} contract, version $1.0.1$, generated from its schema file \texttt{model-metrics.schema.json}: the dotted path of each field down to its leaves (\texttt{[]} an element of a list, \texttt{<key>} any key of a map; a structure met again is listed once, where it first appears), its type, its range or allowed values as the schema states them, whether it is required at its level, its unit and its meaning.}\label{tab:schema-model-metrics}\\
\toprule
Field & Type & Range or values & Req. & Unit & Meaning \\
\midrule
\endfirsthead
\multicolumn{6}{@{}l}{\footnotesize Table~\thetable{} (continued)}\\
\toprule
Field & Type & Range or values & Req. & Unit & Meaning \\
\midrule
\endhead
\bottomrule
\endfoot
\texttt{schema} & const & c\allowbreak{}h\allowbreak{}a\allowbreak{}n\allowbreak{}c\allowbreak{}a\allowbreak{}d\allowbreak{}e\allowbreak{}m\allowbreak{}.\allowbreak{}m\allowbreak{}o\allowbreak{}d\allowbreak{}e\allowbreak{}l\allowbreak{}-\allowbreak{}m\allowbreak{}e\allowbreak{}t\allowbreak{}r\allowbreak{}i\allowbreak{}c\allowbreak{}s\allowbreak{}/\allowbreak{}v\allowbreak{}1 & yes & -- & The document schema identifier, chancadem.model-metrics/v1. \\
\texttt{engine\_\allowbreak{}version} & string & pattern \textasciicircum{}\allowbreak{}\textbackslash{}\allowbreak{}d\allowbreak{}+\allowbreak{}\textbackslash{}\allowbreak{}.\allowbreak{}\textbackslash{}\allowbreak{}d\allowbreak{}\{\allowbreak{}2\allowbreak{}\}\allowbreak{}\textbackslash{}\allowbreak{}.\allowbreak{}\textbackslash{}\allowbreak{}d\allowbreak{}\{\allowbreak{}3\allowbreak{}\}\allowbreak{}\$ & yes & -- & The product version that wrote the document (format X.XX.XXX); a version other than VERSION makes the row stale. \\
\texttt{surrogate} & object & -- & yes & -- & The held-out evaluation of the surrogate against the canonical engine. \\
\texttt{surrogate.\allowbreak{}n\_\allowbreak{}train} & integer & >= 1 & yes & -- & The number of training points of the evaluated bundle. \\
\texttt{surrogate.\allowbreak{}n\_\allowbreak{}test} & integer & >= 1 & yes & -- & The number of points of the independent held-out draw. \\
\texttt{surrogate.\allowbreak{}test\_\allowbreak{}seed} & integer & -- & yes & -- & The seed of the held-out Latin-hypercube draw. \\
\texttt{surrogate.\allowbreak{}per\_\allowbreak{}output} & object & -- & yes & -- & The scores of each output, keyed by its name in the bundle's output\_order. \\
\texttt{surrogate.\allowbreak{}per\_\allowbreak{}output.\allowbreak{}<key>.\allowbreak{}r2} & number & <= 1 & yes & -- & Coefficient of determination of the float32 graph against the engine, in original units. \\
\texttt{surrogate.\allowbreak{}per\_\allowbreak{}output.\allowbreak{}<key>.\allowbreak{}mape\_\allowbreak{}pct} & number & >= 0 & yes & \% & MAPE of the float32 graph against the engine, over the points where the engine's value is not zero. \\
\texttt{surrogate.\allowbreak{}per\_\allowbreak{}output.\allowbreak{}<key>.\allowbreak{}mape\_\allowbreak{}pct\_\allowbreak{}int8} & number & >= 0 & yes & \% & The same MAPE for the int8 graph. \\
\texttt{surrogate.\allowbreak{}p80\_\allowbreak{}monotone\_\allowbreak{}vs\_\allowbreak{}css} & boolean & -- & yes & -- & True when the surrogate's P80 does not fall as the closed-side setting opens on every checked curve. \\
\texttt{surrogate.\allowbreak{}parity} & object & -- & yes & -- & The largest differences between the trained module and its exported graphs, in standardised output units. \\
\texttt{surrogate.\allowbreak{}parity.\allowbreak{}torch\_\allowbreak{}vs\_\allowbreak{}float32\_\allowbreak{}max\_\allowbreak{}abs} & number & >= 0 & yes & -- & The trained torch module against the float32 graph, measured by the train stage right after export. \\
\texttt{surrogate.\allowbreak{}parity.\allowbreak{}float32\_\allowbreak{}vs\_\allowbreak{}int8\_\allowbreak{}max\_\allowbreak{}abs} & number & >= 0 & yes & -- & The float32 graph against the int8 graph through onnxruntime on the held-out draw. \\
\texttt{surrogate.\allowbreak{}note} & string & -- & yes & -- & How the surrogate was evaluated, in words. \\
\texttt{anomaly} & object & -- & yes & -- & The held-out evaluation of the autoencoder beside the engine's envelope check. \\
\texttt{anomaly.\allowbreak{}n\_\allowbreak{}heldout} & integer & >= 1 & yes & -- & The number of held-out points. \\
\texttt{anomaly.\allowbreak{}false\_\allowbreak{}alarm\_\allowbreak{}rate} & number & >= 0; <= 1 & yes & fraction & The share of held-out points the autoencoder marks anomalous. \\
\texttt{anomaly.\allowbreak{}n\_\allowbreak{}planted} & integer & >= 1 & yes & -- & The number of planted off-manifold points with finite features. \\
\texttt{anomaly.\allowbreak{}detection\_\allowbreak{}rate} & number & >= 0; <= 1 & yes & fraction & The share of planted points the autoencoder marks anomalous. \\
\texttt{anomaly.\allowbreak{}envelope\_\allowbreak{}false\_\allowbreak{}alarm\_\allowbreak{}rate} & number & >= 0; <= 1 & yes & fraction & The share of held-out points the envelope check flags. \\
\texttt{anomaly.\allowbreak{}envelope\_\allowbreak{}detection\_\allowbreak{}rate} & number & >= 0; <= 1 & yes & fraction & The share of planted points the envelope check flags. \\
\texttt{anomaly.\allowbreak{}planted} & string & -- & yes & -- & How the off-manifold points were planted (which inputs moved outside the training bounds, and by how much). \\
\texttt{anomaly.\allowbreak{}note} & string & -- & yes & -- & The alarm threshold and how the envelope check is read, in words. \\
\end{longtable}
\endgroup
